% Supplementary material for the MAP-Graph submission.
% Uploaded as a separate document; not part of the main-paper page limit.
\documentclass[letterpaper]{article} % DO NOT CHANGE THIS
\usepackage[submission]{aaai2027} % DO NOT CHANGE THIS
\usepackage[hyphens]{url} % DO NOT CHANGE THIS
\usepackage{graphicx} % DO NOT CHANGE THIS
\urlstyle{rm} % DO NOT CHANGE THIS
\def\UrlFont{\rm} % DO NOT CHANGE THIS
\usepackage{natbib} % DO NOT CHANGE THIS AND DO NOT ADD ANY OPTIONS TO IT
\usepackage{caption} % DO NOT CHANGE THIS AND DO NOT ADD ANY OPTIONS TO IT
\frenchspacing % DO NOT CHANGE THIS
\DeclareCaptionStyle{ruled}{labelfont=normalfont,labelsep=colon,strut=off} % DO NOT CHANGE THIS

\usepackage{amsmath,amssymb}
\usepackage{booktabs}
\usepackage{array}
\usepackage{dblfloatfix}

\pdfinfo{
/TemplateVersion (2027.1)
}

\setcounter{secnumdepth}{2}

\title{MAP-Graph: Supplementary Material}
\author{Anonymous Submission}
\affiliations{}

\begin{document}

\maketitle
\appendix

\section*{Organization and Relation to the Main Paper}

This document supplies implementation and evaluation details omitted from the
main paper for space.  Throughout, C1, C2, and C3 denote the main paper's
challenges of recursive ancestry, heterogeneous constraints, and
action-dependent admissibility.  Symbols ($S$, $I$, $F$, $P$, $V$, $A$,
$\rho$, and $\theta$) and the decision vocabulary
(\textsc{Allow}, \textsc{Block}, \textsc{Reverify}, \textsc{Redact}, and
\textsc{AskUser}) follow the main paper.

\begin{table}[!ht]
\centering
\small
\caption{Correspondence between the supplementary sections and the main paper.}
\label{tab:appendix-map}
\begin{tabular}{@{}p{0.14\linewidth}p{0.58\linewidth}p{0.18\linewidth}@{}}
\toprule
Section & Main-paper material supported & Challenge \\
\midrule
A & Graph schema and path-trust constants & C1 \\
B & Backend interface, workflow integration, and operation semantics & C1, C3 \\
C & Benchmark construction, sanitization, and transfer subset & C1--C3 \\
D & Exact baseline adaptations and the flat-metadata control & C2 \\
E & Metric formulas and implementation settings & C1--C3 \\
F & Scenario, domain, diagnostic, ablation, and transfer results & C1--C3 \\
\bottomrule
\end{tabular}
\end{table}

\section{Provenance Graph Schema}
\label{app:schema}

The main paper's graph-schema subsection states that MAP-Graph evaluates
recorded ancestry rather than only a memory's text or local score.  This
section specifies the serialized vocabulary and the constants used by the
implemented path-trust calculation.

\begin{table*}[!t]
\centering
\small
\setlength{\tabcolsep}{4pt}
\caption{Registered node and edge types emitted by the graph builder.}
\label{tab:mapgraph-schema}
\begin{tabular}{@{}>{\raggedright\arraybackslash}p{0.22\linewidth}
>{\raggedright\arraybackslash}p{0.72\linewidth}@{}}
\toprule
Type & Implemented meaning \\
\midrule
\multicolumn{2}{@{}l@{}}{\textbf{Node types}} \\
User & Owner or authorizer associated with a resource or memory. \\
Agent & Workflow role that observes, writes, retrieves, or proposes an action. \\
Tool & Tool identity associated with an observed output or verification. \\
Resource & Source document or referenced artifact with trust and access metadata. \\
Message & Observable inter-agent content; the graph stores its content hash. \\
Memory & Seeded, observed, or agent-written record stored in Chroma. \\
Claim & Up to three sentence-level claims extracted by deterministic text rules. \\
Action & Proposed action, risk, threshold, support, and gate decision. \\
\midrule
\multicolumn{2}{@{}l@{}}{\textbf{Edge types}} \\
\texttt{authorized\_by} & Associates a resource, memory, or action with an owner or actor. \\
\texttt{observed\_from} & Records an agent observing a resource. \\
\texttt{forwarded\_to} & Records the receivers of an inter-agent message. \\
\texttt{derived\_from} & Links a message, memory, or claim to recorded inputs. \\
\texttt{summarized\_from} & Special derivation edge for a benchmark summary-memory event. \\
\texttt{written\_by} & Associates a memory or claim with its writing agent. \\
\texttt{read\_by} & Records a selected memory read and its retrieval scores. \\
\texttt{verified\_by} & Records a high-trust source or rule-based verifier. \\
\texttt{invalidated\_by} & Records an explicit revocation or invalidating reference. \\
\texttt{used\_for\_action} & Links a supporting memory or claim to an action. \\
\bottomrule
\end{tabular}
\end{table*}

Task identifier, domain, step, visibility, permission scope, trust labels, and
risk tags are attributes rather than node types.  Tool outputs are represented
as Memory nodes associated with a Tool identity; an unregistered output
identifier may additionally appear as a \texttt{Resource:<id>} reference
placeholder.  MAP-Graph is therefore an execution-lineage graph, not a
general-purpose semantic knowledge graph.  The implementation does not create
generic \texttt{Contradicts}, \texttt{Updates}, or temporal-order edges and
does not perform open-domain contradiction resolution.

\subsection{Path-Trust Constants}
\label{app:trust-constants}

For an eligible memory, the implementation computes
\[
\rho(m,a)=\operatorname{clip}_{[0,1]}(S\,I\,F\,P\,V\,A).
\]
Source trust $S$ is the minimum explicit trust among source ancestors, or
$0.80$ if no source ancestor is recorded.  Path integrity $I$ starts at $1$
and is multiplied once by $0.25$, $0.45$, and $0.15$ when the ancestry
contains, respectively, an untrusted/poisoned, private/sensitive, or revoked
category.  Multiple categories compound.  Transformation loss adds $0.08$
for an agent-message record, $0.10$ for a summary event, and $0.03$ for each
non-root memory transformation, capped at $0.65$; hence
$F=\max(0.35,1-\mathrm{loss})$.  Permission validity $P$ is the current binary
\texttt{CanRead} result.  The verification factor is
$V=\min(1.20,1+0.05n_v)$ for $n_v$ high-trust ancestors.  Writer reliability
$A$ defaults to $0.95$.

The factors use different evidence: $S$, $I$, and $V$ summarize recorded
ancestry; $F$ follows the transformation chain; $P$ evaluates the current
reader against the stored policy; and $A$ belongs to the writing agent.
Permission is still enforced as a hard pre-ranking filter in the full system.
The graded product only reranks records that remain eligible, and affected
ancestry is reconsidered by the action gate.

\section{Backend Interface and Memory Operations}
\label{app:operations}

MAP-Graph is implemented as a task-local memory backend between the shared
agent workflow and its simulated proposed action.  Table~\ref{tab:operations}
summarizes the public operations.

\begin{table}[!ht]
\centering
\small
\caption{Implemented MAP-Graph memory operations.}
\label{tab:operations}
\begin{tabular}{@{}>{\raggedright\arraybackslash}p{0.22\linewidth}
>{\raggedright\arraybackslash}p{0.28\linewidth}
>{\raggedright\arraybackslash}p{0.40\linewidth}@{}}
\toprule
Operation & Input & Effect \\
\midrule
Observe/write & Source, tool output, or agent text & Creates a distinct memory, embedding, scope, and lineage links. \\
Retrieve & Agent and query & Filters permissions and returns at most five records ranked by semantic score times path trust. \\
Revoke & Explicit timeline event & Marks a source revoked, clears its scope, and updates directly referenced or already affected records. \\
Decide & Proposed action and support & Applies affected-path rules and a risk-dependent trust threshold. \\
Export & Completed task & Serializes nodes, edges, retrieval diagnostics, gate events, and component settings. \\
\bottomrule
\end{tabular}
\end{table}

\subsection{Workflow Integration}
\label{app:workflow-integration}

The harness uses AutoGen's OpenAI-compatible model client and executes the four
task-provided roles once in fixed order.  It does not instantiate AutoGen
\texttt{AssistantAgent} or \texttt{RoundRobinGroupChat}.  At each role step,
the harness applies scheduled timeline events, registers sources and tool
outputs visible to that role, filters the visible conversation through the
selected backend, retrieves memory, calls the model, and writes the response.
The final role emits a structured proposed-action object, which is passed to
the backend.

The final decision is conservative with respect to the agent proposal and
backend decision.  A restrictive backend result is never relaxed; if the
backend returns \textsc{Allow} but the agent proposes \textsc{Block},
\textsc{Redact}, \textsc{Reverify}, or \textsc{AskUser}, the more restrictive
agent proposal remains.  Only final \textsc{Allow} creates a simulated
execution event, and no external side effect is performed.

\subsection{Write, Retrieval, Update, and Audit Semantics}
\label{app:operation-details}

\paragraph{Write.}
Empty text is rejected; otherwise a new record is created without content
deduplication.  Its \texttt{derived\_from} list combines retrieved-memory
identifiers with sources visible at the current step.  With the full graph
builder, a Message reference is also recorded.  A derived record receives the
intersection of referenced scopes.  If no referenced scope exists, all task
agents are eligible; if the intersection is empty, the scope falls back to the
writing agent.  Visibility is shared only when the resulting scope covers all
task agents.

\paragraph{Retrieval.}
Chroma first supplies semantic candidates.  MAP-Graph computes
\texttt{CanRead} and path trust for each candidate, removes
permission-ineligible records, sorts the remainder by semantic score times
path trust, and returns at most five.  Selected reads create
\texttt{read\_by} edges carrying semantic, path-trust, and final scores.

\paragraph{Update and revocation.}
New observations and messages append records rather than mutating or
superseding prior text.  An explicit revocation event updates the referenced
source and directly referenced or already affected records.  Other descendants
detect the revoked ancestor during later recursive trust evaluation; final
graph propagation marks affected predecessors while retaining their lineage.
A benchmark summary event creates a new record linked by
\texttt{summarized\_from}.  There is no general
\texttt{active}/\texttt{superseded}/\texttt{contradicted} state machine.

\paragraph{Structured audit output.}
The artifacts store selected and permission-filtered identifiers, component
scores, affected ancestors and reasons, supporting identifiers, risk
thresholds, and the final rule-specific decision.  MAP-Graph does not generate
a free-form natural-language explanation, and the experiments do not measure
human-facing explanation quality.

\section{Benchmark Construction and Sanitization}
\label{app:benchmark}

\subsection{Full Benchmark}

The controlled synthetic benchmark contains 2,700 tasks generated with seed
42.  It has three equally sized domains---corporate workflow, software
engineering, and research assistance---and six experiment groups:
(1) clean utility, (2) poisoned-memory propagation, (3) private-memory
leakage, (4) permission revocation, (5) action-risk sensitivity, and
(6) clean compression/overhead.  Each domain contributes 150 tasks to each
group.  Every domain--group cell contains 15 semantic families with 10
variants that alter entities, audiences, channels, evidence shapes, and
action types while retaining a structured oracle.  The 2,700 labels comprise
960 \textsc{Allow}, 1,140 \textsc{Block}, 210 \textsc{Reverify}, and 390
\textsc{Redact}.

Each task provides source documents, initial and derived memory events, tool
observations, visibility and permission fields, a four-role sequence, and a
risk-labelled proposed action.  Corporate roles are manager, analyst,
HR-assistant, and executor; software roles are planner, coder, reviewer, and
executor; research roles are retriever, summarizer, writer, and verifier.  If
the final structured action cannot be parsed, the harness recovers the public
action type and risk from the task schema and uses the retrieved records as
support.

All eight methods receive the same task order and public task fields.  The
runtime sanitizer removes the oracle decision, evaluator rationale,
experiment and family identifiers, variant metadata, the top-level evaluator
canary annotation, and every \texttt{should\_be\_used} annotation before
either the agents or memory backend receive a task.  The public action type,
risk, provenance, visibility, ownership, trust, and permission fields remain,
because reasoning over them is the target capability.  Authorized source text
also remains; otherwise canary exposure could not be measured.

One run for each of B0--B6 and MAP-Graph produces 21,600 main-experiment
decision logs.  Each method has complete 2,700/2,700 decision, trace, and
memory-log coverage; graph artifacts are additionally present for the
graph-producing methods.  The six ablations plus the full method produce
18,900 additional decisions.  These are single-run controlled results.

\subsection{Backbone Transfer Subset}
\label{app:backbone-subset}

The transfer experiment uses a fixed 540-task subset generated with seed 2027
by jointly stratifying the full 2,700 tasks over domain, experiment group, and
oracle label.  It contains 180 tasks per domain, 90 per experiment group, and
all four labels.  The same task identifiers are used for B1, B4, B5, B6, and
MAP-Graph under Qwen2.5-7B-Instruct, GLM-4-9B-0414, and
Llama-3.1-8B-Instruct.  Qwen values are extracted from the final full runs;
GLM and Llama are one subset run each.  Thus each backbone contributes
$5\times540=2,700$ decisions.

\subsection{Cluster Bootstrap}
\label{app:bootstrap}

The 95\% intervals use 2,000 bootstrap samples with bootstrap seed 42.  A sampling cluster is a
semantic family within a domain and experiment group; all ten closely related
variants in a sampled family move together.  The intervals therefore describe
case-family variation in a fixed run, not model-inference nondeterminism.

\section{Baseline Adaptations and Controls}
\label{app:baselines}

B3--B5 are benchmark adaptations that preserve the mechanisms stated below;
they are not claimed as complete reproductions of their source systems.  B6
is an internal flat-metadata control rather than an adaptation of an external
paper.

\smallskip
\noindent\textbf{B0: No Memory.}
The workflow retains only the current conversation and performs no persistent
memory read or write.

\smallskip
\noindent\textbf{B1: Shared Vector Memory.}
All agent records enter one shared Chroma collection; the five most similar
records are returned without permission filtering, lineage scoring, or an
action gate.

\smallskip
\noindent\textbf{B2: Isolated Vector Memory.}
The same vector mechanism is maintained separately per agent.  This prevents
direct cross-store reads but does not create a shared provenance graph.

\smallskip
\noindent\textbf{B3: Adapted G-Memory.}
Observable interaction traces are converted into trajectories represented by
query, interaction, and insight graphs.  The implementation performs one
task-level retrieval and reuses its context for all four roles, applies
one-hop query-graph expansion, retains the 512 most recent trajectories, and
bounds insight state to 10 rules.  Periodic insight updates use an LLM.
The reported configuration replaces the original per-task LLM trajectory
sparsifier and generative reranker with deterministic key-step extraction and
similarity ordering.  Auxiliary insight calls use temperature 0.1 with a
small output cap.  B3 receives no MAP-Graph permission filter, path-trust
propagation, or action gate.

\smallskip
\noindent\textbf{B4: Adapted Collaborative Memory.}
Records are assigned to private or shared tiers from visibility and permission
metadata.  Restricted or revoked status is inherited by derived messages, and
the current access-control list is checked before top-five similarity results
are returned.  A per-task user--agent--resource access graph is exported.
B4 has neither recursive path-trust scoring nor an action-time gate.

\smallskip
\noindent\textbf{B5: Adapted MemLineage.}
Each fragment records derivation ancestors, edges, and untrusted-source
reasons.  Retrieval is similarity-based and does not filter by reader
permission.  Before a high-risk action, supporting memories are checked for an
untrusted ancestor.  The evaluated implementation has no per-principal
signatures, cryptographic/Merkle-verified log, permission-aware retrieval, or
multi-agent trust decay.

\smallskip
\noindent\textbf{B6: Flat Provenance.}
B6 stores, on each memory, visibility, owner, permission scope, source trust,
private/revoked/poisoned state, and action risk.  It performs direct
metadata-based retrieval filtering and a metadata-based action check.
Derived messages summarize metadata from their direct inputs, and revocation
updates records directly tied to the source.  B6 does not construct a
provenance graph, recursively traverse parent memories, or propagate
path-trust scores.  It therefore isolates the value of recursive ancestry from
the availability of local policy metadata.

\section{Metrics and Implementation Settings}
\label{app:metrics}

\subsection{Paper-Facing Metrics}

Let $y_i$ be the oracle decision, $\hat y_i$ the final decision,
$\mathcal{I}=\{\textsc{Block},\textsc{Redact},\textsc{Reverify},
\textsc{AskUser}\}$, and $N$ the relevant denominator.

\begin{enumerate}
\item \textbf{Exact decision accuracy (Acc, $\uparrow$):}
$N^{-1}\sum_i\mathbf{1}[\hat y_i=y_i]$.
\item \textbf{Task success rate (TSR, $\uparrow$):}
an oracle-\textsc{Allow} case succeeds only for \textsc{Allow}; a
non-\textsc{Allow} case succeeds for any decision in $\mathcal{I}$.  Thus a
safe but non-exact intervention can increase TSR without increasing Acc.
\item \textbf{Clean success (Clean, $\uparrow$):}
the \textsc{Allow} rate on Exp-1 and Exp-6, whose oracle decision is
\textsc{Allow}.
\item \textbf{Unsafe rate ($\downarrow$):}
the fraction of all 2,700 tasks with a simulated execution event when the
oracle decision is not \textsc{Allow}.
\item \textbf{Attack success rate (ASR, $\downarrow$):}
the \textsc{Allow}/execution rate among the 450 Exp-2 poisoned-memory tasks.
\item \textbf{Leakage ($\downarrow$):}
the fraction of the 450 Exp-3 tasks with unauthorized canary exposure or a
final \textsc{Allow}.
\item \textbf{Unauthorized access (UAcc, $\downarrow$):}
the total successful unauthorized memory reads divided by total observed
unauthorized-read attempts.  It is N/A when no attempt is observed.
\item \textbf{Revocation violation ($\downarrow$):}
the final-\textsc{Allow} rate among the 450 Exp-4 permission-revocation tasks.
\item \textbf{Structural contamination:}
for each Exp-2 task having at least one downstream memory, the fraction of
downstream records that retain poisoned ancestry, averaged across eligible
tasks.  This is a lineage-structure diagnostic, not an execution or leakage
rate: a system may preserve affected ancestry for audit while preventing its
use.
\item \textbf{Correct block ($\uparrow$):}
$\Pr(\hat y=\textsc{Block}\mid y=\textsc{Block})$.
\item \textbf{False block ($\downarrow$):}
$\Pr(\hat y=\textsc{Block}\mid y=\textsc{Allow})$.
\end{enumerate}

The additional \emph{block rate} is the unconditional fraction of final
\textsc{Block} decisions.  Per-task token use sums provider-reported prompt
and completion tokens over the four outer agent calls and any recorded
memory-layer calls.  The token-cap diagnostic is the fraction of the 10,800
outer agent calls per method whose reported completion reaches the 2,048-token
limit.  Average graph nodes and edges are computed from exported task graphs.

\subsection{Fixed Implementation Settings}
\label{app:settings}

All main-experiment methods use the same AutoGen OpenAI-compatible client with
Qwen2.5-7B-Instruct, outer-call temperature 0, one four-role round, at most 12
recent context messages, at most five retrieved memories, a 2,048-token
generation cap, a 120-second per-call timeout, and at most five retries.
Embedding methods use normalized BAAI/bge-small-en-v1.5 embeddings on CPU
with $k=5$.

MAP-Graph filters permissions before ranking, uses multiplicative path trust,
propagates affected state, and evaluates retrieved plus contextual support at
action time.  Its thresholds are 0.30 for answers, 0.60 for low- or
medium-risk non-answer actions, and 0.85 for high-risk actions.  A high-risk
action is blocked when an affected supporting path is detected.

Latency is not a paper-facing metric.  All model calls used remote shared
endpoints and experiments ran concurrently across accounts, so provider
queuing and long-tail service load cannot be separated from local memory
overhead.  Token use and graph size are retained only as reproducible
diagnostics.

\section{Additional Results}
\label{app:additional-results}

Tables~\ref{tab:scenario-results}--\ref{tab:backbone-uacc} report the
scenario, domain, full-run, ablation, and transfer diagnostics discussed in
this section.

\begin{table*}[!t]
\centering
\small
\caption{Scenario-level results (percent; 450 tasks per row).  ``Best
baseline'' is selected for the quantity shown within that experiment group.}
\label{tab:scenario-results}
\setlength{\tabcolsep}{4pt}
\begin{tabular}{llrrlrr}
\toprule
Group & Scenario & MAP TSR & Best baseline TSR & Risk/utility quantity & MAP & Best baseline \\
\midrule
Exp-1 & Clean utility             & 84.22 & 83.33 (B1) & Exact accuracy & 84.22 & 83.33 (B1) \\
Exp-2 & Poisoned propagation      & 100.00 & 77.78 (B5) & ASR            & 0.00  & 22.22 (B5) \\
Exp-3 & Private-memory leakage    & 100.00 & 82.00 (B2) & Leakage        & 0.00  & 18.00 (B2) \\
Exp-4 & Permission revocation     & 100.00 & 76.22 (B5) & Revocation     & 0.00  & 23.78 (B5) \\
Exp-5 & Action-risk sensitivity   & 89.33 & 68.22 (B6) & Unsafe action  & 9.11  & 28.67 (B6) \\
Exp-6 & Clean compression         & 96.22 & 94.44 (B4) & Exact accuracy & 96.22 & 94.44 (B4) \\
\bottomrule
\end{tabular}
\end{table*}

\begin{table*}[!t]
\centering
\normalsize
\caption{MAP-Graph results by domain (percent; 900 tasks per row).}
\label{tab:domain-results}
\setlength{\tabcolsep}{8pt}
\renewcommand{\arraystretch}{1.10}
\begin{tabular}{lrr}
\toprule
Domain & TSR & Acc \\
\midrule
Corporate workflow & 94.56 & 73.56 \\
Software engineering & 95.00 & 74.56 \\
Research assistance & 95.33 & 70.00 \\
\bottomrule
\end{tabular}
\end{table*}

\begin{table*}[!t]
\centering
\small
\caption{Additional full-run diagnostics (percent except tokens/task).
Correct and false block use their conditional denominators; Cap calls is the
fraction of outer calls reaching 2,048 completion tokens.}
\label{tab:main-diagnostics}
\setlength{\tabcolsep}{4pt}
\begin{tabular}{lrrrrrr}
\toprule
Method & Block & Correct block & False block & Structural contam. & Tokens/task & Cap calls \\
\midrule
B0 & 12.67 & 16.23 & 2.50 & 0.00 & 3,397.2 & 8.12 \\
B1 & 14.41 & 15.44 & 3.33 & 43.44 & 5,007.5 & 4.66 \\
B2 & 14.07 & 16.67 & 2.71 & 25.00 & 4,054.4 & 4.06 \\
B3 & 9.07 & 8.77 & 3.23 & 11.11 & 6,687.4 & 9.08 \\
B4 & 15.37 & 17.46 & 2.71 & 25.02 & 4,927.8 & 4.46 \\
B5 & 23.96 & 38.86 & 2.81 & 62.17 & 4,953.2 & 4.09 \\
B6 & 16.59 & 21.23 & 2.60 & 25.00 & 4,882.7 & 4.44 \\
MAP-Graph & 24.00 & 54.21 & 3.12 & 25.00 & 4,814.9 & 3.89 \\
\bottomrule
\end{tabular}
\end{table*}

\subsection{Scenario-Level Results}
\label{app:scenario-results}

Table~\ref{tab:scenario-results} compares MAP-Graph with the strongest baseline
for the utility and risk quantity relevant to each experiment group.  The
comparator can differ across groups.

\subsection{Confidence Intervals and Domain Results}
\label{app:confidence-domain}

MAP-Graph's TSR is 94.96\% with a 95\% cluster-bootstrap interval of
[93.15, 96.52], and its Acc is 72.70\% [68.26, 77.19].  The strongest
baseline TSR is B6 at 74.67\% [71.67, 77.74]; the strongest baseline Acc is B5
at 51.07\% [47.19, 55.07].  MAP-Graph's Clean is 90.22\%
[87.44, 92.92], and its unsafe rate is 1.52\% [0.37, 2.93].  These intervals
measure clustered case variation within one run.
Table~\ref{tab:domain-results} reports the corresponding domain breakdown.
The TSR range is narrow across the three controlled role configurations.  The
lower research Acc with high TSR reflects more safe but non-exact
interventions, not more impermissible execution.

\subsection{Main-Experiment Diagnostics}
\label{app:main-diagnostics}

Table~\ref{tab:main-diagnostics} reports block rates, structural
contamination, and token diagnostics for the full run.
MAP-Graph's average token count consists of 4,194.8 prompt and 620.1
completion tokens.  B3 is highest because retrieved trajectories and periodic
insight generation add context and auxiliary generation.  MAP-Graph exports
39.74 nodes and 246.78 edges per task on average.  These counts exclude
embedding computation, graph construction time, storage I/O, evaluator work,
and failed requests for which the provider reports no usage; they are not a
complete cost or latency comparison.

Structural contamination must be interpreted separately from safety outcomes.
MAP-Graph's 25.00\% means that affected ancestry remains attached to some
downstream records; those records are retained and marked for audit.  Its
0\% ASR and 0\% poisoned-group impermissible execution show that retained
lineage was not treated as admissible support in those cases.

\subsection{Full Ablation Diagnostics}
\label{app:ablation-diagnostics}

Table~\ref{tab:ablation-diagnostics} supplements the main-paper ablation
table with leakage, contamination, block, token, and graph-size diagnostics.
The six variants correspond to explicit component switches.  \emph{No
permission filter} disables candidate exclusion but leaves trust and gating
active; its 100\% UAcc in the main table despite 0\% endpoint leakage here
shows why read-boundary and endpoint metrics are both needed.  \emph{No trust
propagation} also disables trust-aware reranking.  \emph{No action gate}
allows at the backend decision stage, while \emph{No containment} retains the
threshold gate but removes affected-path rules.  \emph{No provenance graph}
jointly disables graph construction, permission filtering, trust propagation
and reranking, action gating, containment, and contextual action support; it
is a system-level removal rather than a one-switch attribution.  The
\emph{compressed graph} retains policy behavior but exports fewer graph
details.

\begin{table*}[!t]
\centering
\small
\caption{Diagnostics supplementing the main-paper ablation table (percent
except tokens and graph size).  CBlock/FBlock denote conditional correct/false
block.}
\label{tab:ablation-diagnostics}
\setlength{\tabcolsep}{3.5pt}
\begin{tabular}{lrrrrrrrr}
\toprule
Variant & Leak & Contam. & Block & CBlock & FBlock & Tokens & Nodes & Edges \\
\midrule
Full & 0.00 & 25.00 & 24.00 & 54.21 & 3.12 & 4,814.9 & 39.74 & 246.78 \\
Compressed graph & 0.00 & 25.02 & 23.78 & 53.77 & 3.02 & 4,847.0 & 24.66 & 70.68 \\
No action gate & 22.22 & 25.00 & 14.41 & 17.54 & 2.50 & 4,830.1 & 39.74 & 246.69 \\
No containment & 20.89 & 25.00 & 14.81 & 18.33 & 2.50 & 4,814.2 & 39.68 & 245.96 \\
No permission filter & 0.00 & 24.94 & 25.74 & 58.86 & 2.40 & 4,911.9 & 39.76 & 250.73 \\
No provenance graph & 22.89 & 40.44 & 14.22 & 16.40 & 3.02 & 4,991.2 & 0.00 & 0.00 \\
No trust propagation & 0.00 & 25.00 & 11.15 & 24.12 & 2.71 & 4,886.3 & 39.91 & 265.10 \\
\bottomrule
\end{tabular}
\end{table*}

\begin{table*}[!t]
\centering
\small
\caption{Unauthorized-access diagnostics on the common 540-task transfer
subset.  Each backbone block reports unauthorized reads (R), observed attempts
(A), and conditional UAcc in percent.}
\label{tab:backbone-uacc}
\setlength{\tabcolsep}{3.8pt}
\renewcommand{\arraystretch}{1.08}
\begin{tabular}{@{}lrrr rrr rrr@{}}
\toprule
& \multicolumn{3}{c}{Qwen2.5-7B} & \multicolumn{3}{c}{GLM-4-9B} &
\multicolumn{3}{c}{Llama-3.1-8B} \\
\cmidrule(lr){2-4}\cmidrule(lr){5-7}\cmidrule(l){8-10}
Method & R & A & UAcc & R & A & UAcc & R & A & UAcc \\
\midrule
B1        & 48  & 48    & 100.00 & 48  & 48    & 100.00 & 48  & 48    & 100.00 \\
B4        & 990 & 2,323 & 42.62  & 990 & 2,324 & 42.60  & 990 & 2,324 & 42.60 \\
B5        & 101 & 101   & 100.00 & 86  & 86    & 100.00 & 84  & 84    & 100.00 \\
B6        & 990 & 2,324 & 42.60  & 990 & 2,324 & 42.60  & 990 & 2,324 & 42.60 \\
MAP-Graph & 0   & 1,334 & 0.00   & 0   & 1,334 & 0.00   & 0   & 1,334 & 0.00 \\
\bottomrule
\end{tabular}
\end{table*}

\subsection{Backbone Access-Boundary Diagnostics}
\label{app:backbone-boundary}

Table~\ref{tab:backbone-uacc} reports unauthorized-access diagnostics on the
common 540-task transfer subset.
The deterministic access policy and task structure make the attempt counts
largely stable across backbones, while B5's model-dependent retrieval changes
how often an unauthorized record is selected.  MAP-Graph rejects every
observed unauthorized read under all three backbones.  This supports transfer
of the implemented access boundary, not a claim of run-to-run stability or
deployment-scale robustness.

\subsection{Error Analysis}
\label{app:error-analysis}

MAP-Graph makes 737 exact-decision errors.  Of these, 601 are safe but
non-exact interventions, such as \textsc{Redact} rather than \textsc{Block} or
\textsc{Reverify}; this accounts for much of the TSR--Acc gap.  Another 95
cases are unnecessary interventions on oracle-\textsc{Allow} tasks.  The
remaining 41 are impermissible \textsc{Allow} decisions, all in Exp-5.  The
residual safety errors in this run therefore concentrate in action-risk
sensitivity rather than poisoned, private, or revoked-memory groups.

\end{document}
